\section{STM32CubeIDE Configuration}
\label{chap:test_enviroment}

As mentioned in the section~\ref{chap:Back_ARM_STM32CubeIDE}, all firmware development for 
this project was performed using \texttt{STM32CubeIDE} and its associated tools. This 
section describes the pre-configuration carried out in \texttt{STM32CubeIDE} prior to 
defining the functional tests.

\texttt{STM32CubeIDE} allows the creation of projects specifically targeted to 
particular microcontrollers or development boards. Using this feature, a new project 
is created using a preconfigured STM32F4-based evaluation board as the base. This board 
uses the STM32F407 microcontroller, which, as previously noted, is compatible with 
STM32F405 configurations.

Once the project is created, \texttt{STM32CubeMX} opens automatically, allowing the 
microcontroller to be configured visually as follows:

\begin{enumerate}
  \item Enable I2C1 and configure it for 100~kHz operation with 7-bit addressing. 
  Then assign the I2C1\_SCL and I2C1\_SDA lines to pins PB6 and PB7, respectively.
  
  \item Configure a UART peripheral to enable debugging.
  
  \item Reconfigure pin PA1, which is wired to the LIS3DH INT1 pin, to 
  \textbf{GPIO\_EXTI1} mode. In the configuration for this pin, select 
  ``External interrupt mode with rising edge trigger detection'' and ``No pull-up and 
  no pull-down''. Then, in the \textbf{NVIC} tab, enable the \textbf{EXTI line 1 
  interrupt} and assign it a priority of 5.
\end{enumerate}

Once these configurations are applied and saved, \texttt{CubeMX} generates the code 
required to initialize the I$^2$C, GPIO, and EXTI lines. Within this generated code, and 
before integrating the LIS3DH library or defining the tests, the following two 
modifications must be made:

\begin{enumerate}
  \setcounter{enumi}{3}
  \item In the \texttt{main()} function, locate and comment out the 
  \texttt{SystemClock\_Config()} call.
  
  \item Define the \texttt{HAL\_GPIO\_EXTI\_Callback()} function, which is executed 
  automatically when the DRDY interrupt is triggered on INT1.
\end{enumerate}

The \texttt{SystemClock\_Config()} function is commented out because the STM32F405 
model implemented in QEMU does not accurately emulate the entire clock subsystem. As 
a result, some function calls within \texttt{SystemClock\_Config()} return 
\texttt{HAL\_ERROR}, which in turn causes \texttt{Error\_Handler()} to be invoked and 
enter its infinite loop.
